\section{logic}

If, as Galileo stated, mathematics is the language of nature, it is important
that mathematics itself is expressed in a language that is as objective and
unambiguous as possible. \emph{Logic} is the mathematical discipline that deals
with the study and formalization of mathematical language. In this chapter, we
summarize in a concise and intuitive manner some concepts and simultaneously
establish the notations that will be used subsequently.

Logic studies \emph{formal systems}%
\mymargin{formal systems}%
\index{formal systems} (also known as \emph{logical-deductive systems} or \emph{axiomatic systems}) which are mechanical, unambiguous descriptions of a formal language. We speak in the plural of \emph{formal systems} because each field of mathematics (or other sciences) could develop its own specialized formal system for that field. The first formal system was developed by Euclid to describe the properties of points, lines, and circles in the plane (the \emph{Euclidean geometry}). At Euclid's time, the formalization was still intuitive and incomplete; the modern formalization of this system was completed by Hilbert, two thousand years later. Peano introduced an axiomatic system to describe the properties of natural numbers. Dedekind identified the axioms to describe real numbers. Cantor introduced the \emph{set theory} within which it was possible to include all other mathematical theories. This theory was formalized by Zermelo and Fraenkel, and it is this theory that we will use in our course.

All formal systems describe a language. To describe a language, we must first
specify what the \emph{symbols}%
\mymargin{symbols}%
\index{symbols} of that language are. In general, we can think of symbols as letters (or characters, in computer language) that can be used to compose sentences. Typical symbols of mathematical theories are, for example: \texttt{x}, \texttt{y}, \texttt{5}, \texttt{7}, \texttt{+}, \texttt{=}, \texttt{)}, \texttt{:}, etc. In computer languages, symbols correspond to the characters present on a computer keyboard. In natural languages, one would take the individual letters of the alphabet as symbols, possibly adding punctuation characters. A finite sequence of symbols is called a \emph{formula}%
\mymargin{formula}%
\index{formula} (we could also call them \emph{sentences}). For example: \texttt{x7)+} could be a formula consisting of four symbols. Among all formulas, a formal system must identify those that could be called \emph{well-formed formulas}, i.e., the formulas to which we actually want to assign meaning. For example, the formula $x+5=7$ could be a well-formed formula because we are able to assign it a meaning. The formal system does not assign meaning to well-formed formulas (meaning is an extrapolation of our mind) but simply must provide rules to determine which are well-formed formulas and which are not. Since every symbol we use can be represented on a computer, we can think of formulas as strings in programming languages, and we can think that the formal system must describe an algorithm (i.e., a mechanical procedure) capable of determining whether a formula is well-formed or not. Among the \emph{well-formed formulas}, the formal system must finally specify which are the \emph{theorems}. Again, this must be done through a purely mechanical algorithm to ensure that theorems are objective and universal: there can be no disagreement on the validity of a theorem; there may be disagreement on the interpretation of such a theorem. Typically, formal systems are \emph{deductive}. In \emph{deductive} systems, some formulas are identified and called \emph{axioms}, which are immediately recognized as theorems. For example, we will see that the first axiom of set theory is 
\[
  \exists X\colon \not \exists y\colon y\in X
\]
which could be read as 
\begin{displayquote}
there exists a set $X$ for which no $y$ is an element of $X$
\end{displayquote}
and expresses the existence of the empty set. In addition to axioms, a deductive
system specifies \emph{inference rules}, which are ways in which formulas can be
modified or composed such that if the starting formulas are theorems, the
resulting formula is also a theorem. Aristotle's \emph{syllogisms} can be used
as examples of inference rules. Suppose the formulas \texttt{"Socrates is a
man"} and \texttt{"man is an animal"} are both theorems. Then we can define an
inference rule that tells us that if \texttt{"X is a Y"} is a theorem and
\texttt{"Y is a Z"} is a theorem, then \texttt{"X is a Z"} is also a theorem.
With this inference rule, it is possible to deduce that \texttt{"Socrates is an
animal"} is a theorem. It is clear that if formal rules can be mechanically
defined through an algorithm, then theorems can also be mechanically determined
through an algorithm. Research in mathematics consists of exploring the space of
well-formed formulas to determine which are actually theorems. Providing the
\emph{proof}%
\mymargin{proof}%
\index{proof} of a theorem means exhibiting the entire chain of formulas and inference rules that allow obtaining the theorem from the axioms.

\subsection{propositions, logical operators}

All this is a mechanical process, but in reality, it is obvious that formal
systems are defined in such a way as to have some kind of intuitive meaning for
us. In this way, the search for proofs is not a purely mechanical process but
follows lines of thought that may require intuition, inventiveness, and even
aesthetic sense. Typically, at an intuitive level, we want to assign a truth
value to well-formed formulas: we would like to say that some formulas are
\emph{true}%
\mymargin{true}%
\index{true} and others are \emph{false}%
\mymargin{false}%
\index{false}. In such a case, well-formed formulas are usually called \emph{propositions}%
\mymargin{propositions}%
\index{proposition} (in natural language, we would say: \emph{statements}). An example of a (false) proposition could be: \texttt{2+2=5}.

It is possible to combine multiple propositions using logical operators. If $P$
and $Q$ are propositions, one can construct the proposition $P \land Q$, called
\emph{logical conjunction}%
\mymargin{conjunction}%
\index{logical conjunction}. This proposition can be read as ``$P$ and $Q$'' and is a proposition that is true only if both $P$ and $Q$ are true (see table~\ref{tab:verita_operatori_logici} for a schematic summary). Often, logical conjunction is implied: if a list of propositions $P,Q,R$ is made, it is usually understood as their conjunction $P \land Q \land R$, meaning they must all be true. The \emph{logical disjunction}%
\mymargin{disjunction}%
\index{logical disjunction} denoted by $P \lor Q$ can be read as ``$P$ or $Q$'' and is a proposition that is true if at least one of $P$ or $Q$ is true. The \emph{logical negation}%
\mymargin{negation}%
\index{logical negation} denoted by $\lnot P$ is a proposition that can be read as ``not $P$'' and is true when $P$ is false and false when $P$ is true.

Logical operators that are frequently used are \emph{implications}%
\mymargin{implications}%
\index{implication}. The proposition $P\Rightarrow Q$ can be read as ``$P$ implies $Q$'' and means that $Q$ is true if $P$ is true. Do not confuse the truth value of $P\Rightarrow Q$ with the truth value of $Q$. If $P$ is true, then $P\Rightarrow Q$ is true or false depending on whether $Q$ is true or false. But if $P$ is false, then the implication $P\Rightarrow Q$ is true regardless of the value of $Q$. In fact, $P\Rightarrow Q$ is equivalent to $Q \lor \lnot P$ because for the truth of $P\Rightarrow Q$, it suffices that $Q$ is true (when $P$ is true) or that $P$ is false.

The reverse arrow $P\Leftarrow Q$ can be used to invert the implication: it is
equivalent to $Q \Rightarrow P$. If both implications $(P \Leftarrow Q) \land
(P\Rightarrow Q)$ hold, it is easy to convince oneself that $P$ and $Q$ must
have the same truth value: we will then say that they are equivalent and write
$P \Leftrightarrow Q$.

In table~\ref{tab:verita_operatori_logici}, all the truth values that can be
obtained by combining two propositions are reported. In
table~\ref{tab:operatori_logici}, some properties of these operators are
reported: these properties can be understood by interpreting their meaning but
can also be mechanically deduced. To mechanically verify the validity of one of
these logical expressions, it is sufficient to create a table in which all
possible truth values of $P$, $Q$, and $R$ are inserted (in total, there will be
$8$ cases), and for each of these, the truth value of each operation performed
must be calculated, verifying that each complete expression assumes the value
$T$ (true). We will say that these are \emph{tautologies}%
\index{tautology}%
\mymargin{tautology}%
i.e., logical expressions whose validity does not depend on the truth value of
its terms.

\begin{table}
\begin{center}
  \begin{tabular}{cc|cccccc}
    $P$ & $Q$ & $\neg P$ & $P\land Q$ & $P\lor Q$ & $P\Rightarrow Q$ &
    $P\Leftarrow Q$ & $P\Leftrightarrow Q$ \\\hline
    \texttt{F} & \texttt{F} & \texttt{V} & \texttt{F} & \texttt{F} & \texttt{V} & \texttt{V} & \texttt{V} \\
    \texttt{F} & \texttt{V} & \texttt{V} & \texttt{F} & \texttt{V} & \texttt{V} & \texttt{F} & \texttt{F} \\
    \texttt{V} & \texttt{F} & \texttt{F} & \texttt{F} & \texttt{V} & \texttt{F} & \texttt{V} & \texttt{F} \\
    \texttt{V} & \texttt{V} & \texttt{F} & \texttt{V} & \texttt{V} & \texttt{V} & \texttt{V} & \texttt{V} \\
    \end{tabular}
\end{center}
\caption{The truth table of logical operators. 
$\texttt{F}$ means \emph{false}, $\texttt{V}$ means \emph{true}.}
\label{tab:verita_operatori_logici}
\end{table}

\begin{table}
\begin{tabular}{rcll}
                                                    &$\neg (P \land \neg P)$&
                          & non-contradiction \\
                                                  &$P \lor \neg P$&
                         & excluded middle \\
                                                  $\neg \neg P$ & $\iff$ & $ P$
                         & double negation\\
                                                                        $P \land
                                    Q$ & $\iff$ & $ Q \land P$
                                    & symmetry\\
                                                                          $P
                                     \lor Q$ & $\iff$ & $ Q \lor P$
                                     & \\
                                                            $\neg (P\land Q)$ &
                              $\iff$ & $ (\neg P) \lor (\neg Q)$      & De
                              Morgan's laws\\
                                                              $\neg (P\lor Q)$ &
                               $\iff$ & $ (\neg P) \land (\neg Q)$     & \\
                                                        $(P\land Q) \lor R$ &
                            $\iff$ & $ (P\lor R) \land (Q \lor R)$  &
                            distributive property\\
                                                        $(P\lor Q) \land R$ &
                            $\iff$ & $ (P\land R) \lor (Q \land R)$ & \\
                                                        $(P \Rightarrow Q)$ &
                            $\iff$ & $ (Q \Leftarrow P)$            &
                            antisymmetry\\
                                                        $(P\Rightarrow Q)$ &
                            $\iff$ & $ (\neg Q\Rightarrow\neg P)$   &
                            contrapositive implication\\
                                                $\neg (P\Rightarrow Q)$ & $\iff$
                        & $ P \land (\neg Q)$            & counterexample\\
                                                          $P\Rightarrow Q$ &
                             $\iff$ & $ \lnot(P \land (\neg Q))$     & proof by
                             contradiction
\end{tabular}
\caption{Some properties of logical operators. These propositions are all true regardless of the truth values of the propositions $P$ and $Q$. In particular, logical equivalences tell us that the propositions on both sides of the equivalence always assume the same truth value and are therefore interchangeable.}
\label{tab:operatori_logici}
\end{table}

To formally define the logic of propositional calculus, we must have an
underlying system that tells us what the basic propositions are, i.e., the terms
we can substitute in place of $P$, $Q$, $R$. These will be propositions (i.e.,
well-formed formulas). Once the basic propositions are defined, any formula of
the form $(P)\land(Q)$, $(P)\lor(Q)$, $(P)\implies(Q)$, $(P)\impliedby (Q)$,
$(P)\Leftrightarrow(Q)$, and $\lnot(P)$ is also a proposition (i.e.,
well-formed) if $P$ and $Q$ are. The formal system uses the symbols: $\land$
$\lor$ $\implies$ $\impliedby$ $\Leftrightarrow$ $\lnot$ $($ $)$ in addition to
all the symbols needed to compose the basic propositions.

The simplest possible case is one in which we have only two basic propositions
$F$ and $V$ and two axioms: $V$ and $\lnot (F)$. The interpretation of the
system is that $V$ is true while $F$ is false.

The inference rules of propositional calculus can be given in different ways. 
\index{inference rules}%
\mymargin{inference rules}%
Let us try to explain (without too many formalities) the system called
\emph{natural deduction}. This system presents, for each of the logical
operators (negation, conjunction, disjunction, implication), two formal rules:
one that allows the introduction of the operator and one that allows the
elimination of the operator.

The simplest rules are those concerning logical conjunction. If $P$ and $Q$ are
theorems, we can deduce that $(P)\land (Q)$ is also a theorem (introduction of
conjunction).
\index{logical conjunction}%
\index{introduction!conjunction}%
\mymargin{introduction conjunction}%
Conversely, if $(P)\land (Q)$ is a theorem, we can deduce that both $P$ and $Q$
are theorems (elimination of conjunction).
\index{elimination!conjunction}%
\mymargin{elimination conjunction}%
For example, if we have the theorem $1+1=2$ and the theorem $2+1=3$, then we can
deduce that $(1+1=2) \land (2+1=3)$ is a theorem. Conversely, if $(1+1=2)\land
(2+1=3)$ is a theorem, then we can deduce that $2+1=3$ is a theorem.

The introduction of logical disjunction is also a very simple rule: if $P$ is a
theorem, then $(P)\lor (Q)$ and $(Q)\lor (P)$ are theorems, regardless of the
proposition $Q$.
\index{introduction disjunction}%
\mymargin{introduction disjunction}%
For example, if $1+1=2$ is a theorem, we can deduce that $(1+1=2) \lor (1+1\neq
2)$ is a theorem. The rule of elimination of disjunction is slightly more
complicated. If $(P)\lor (Q)$ is a theorem and both $(P)\implies (R)$ and
$(Q)\implies (R)$ are theorems, we can deduce that $R$ is a theorem.
\mymargin{elimination disjunction}%
For example, if we have the theorems: 
\[
\begin{gathered}
  (x>1) \implies (x^2>1)\\ 
  (x<-1) \implies (x^2>1)\\  
  (x>1) \lor (x<-1)
\end{gathered}
\]
then we can deduce that $x^2>1$ is a theorem.

Regarding logical implication, the rule of elimination is called \emph{modus
ponens} and is perhaps the most important of the inference rules.
\index{modus ponens}%
\mymargin{modus ponens}%
The rule tells us that if we have the theorem $P$ and the theorem $(P)\implies
(Q)$, we can assert that $Q$ is also a theorem. For example, if we have proven
$(1+1=2) \implies (2+1=3)$ and we have proven $1+1=2$, then we can deduce
$2+1=3$. This formal rule gives meaning to logical implication.

The rule of \emph{introduction of implication} is more complicated to formalize
and requires the introduction of a new concept: the \emph{demonstrative
context}.
\index{demonstrative context}%
\mymargin{demonstrative context}%
A \emph{context} is essentially a new formal system that inherits all the axioms
and formal rules of the original formal system but with additional properties
that typically allow obtaining theorems that were not possible in the original
formal system. For the introduction of logical implication, the rule is as
follows. To assert that $(P)\implies (Q)$ is a theorem, we introduce a new
context in which all previously proven theorems are axioms, and additionally,
$P$ is an additional axiom. If in this context we manage to obtain $Q$ as a
theorem, then we can say that $(P)\implies (Q)$ is a theorem in the original
formal system. Note that $Q$ might not be a theorem in the original context; it
is only in the hypothetical context where we have assumed (hypothetically) that
$P$ was a theorem.

Let us try, for example, to prove the most controversial rule of the truth
tables, namely that $F\implies V$ (where $F$ stands for false and $V$ stands for
true). What we need to do is verify that if $\lnot P$ is a theorem (i.e., $P$ is
false) and if $Q$ is a theorem (i.e., $Q$ is true), then $P\implies Q$ is a
theorem. We should use the rule of introduction of logical implication. So,
suppose hypothetically that $P$ is a theorem. Since $Q$ was a theorem in the
original context, $Q$ remains a theorem in the hypothetical context where we
have assumed that $P$ is an axiom. We can then, outside the context, deduce
$P\implies Q$ as was to be demonstrated.

For example, instead of $P$, we can substitute $2\neq 2$, and for $Q$, we can
substitute $1=1$. We would then want to prove that $2\neq 2 \implies 1=1$ is a
theorem. Let us try to do it formally with all the details.

Suppose we have a formal system containing the symbols $1$, $2$, $=$, $\neq$ in
addition to all the symbols of propositional logic (i.e., $\land$, $\lor$,
$\lnot$, $\implies$, $($ and $)$). We call the symbols $1$ and $2$
\emph{constants}. Suppose that all formulas of the form $x\neq y$ and $x=y$ are
propositions (well-formed formulas) if we substitute any constant for $x$ and
$y$. Suppose that $x\neq y$ is synonymous with $\lnot (x=y)$. Finally, suppose
that $x=x$ is an axiom if we substitute any constant for $x$.

Outside the formal system, when presenting a proof, we will list our theorems,
numbering each line for reference. We use curly braces $\{$ and $\}$ to enclose
demonstrative contexts. In square brackets, we indicate the formal rules we are
using.
\begin{theorem}[example of formal proof]
  $2 \neq 2 \implies 1=1$.
\end{theorem}
%  
\begin{proof}
  \begin{align}
    &\{ & &\text{[begin hypothetical context]} \\ 
    &\quad 2\neq 2 & & \text{[hypothesis]}  \\
    &\quad 1=1 & & \text{[axiom]} \\
    &\} & & \text{[end of hypothetical context]} \\
    &2\neq 2 \implies 1=1 & & \text{[introduction of implication]}
  \end{align}
\end{proof}

% For example, we can prove that 
% $1=2 \implies (2=3 \implies 1=3)$.
% We start by introducing a hypothetical context in which $1=2$.
% In this context, we introduce a new \emph{sub}-context 
% in which $2=3$. In the sub-context, we have as (hypothetical) theorems
% both $1=2$ and $2=3$. 
% An axiom of equality (which we have not yet mentioned)
% will allow us to say that if $x=y$, then in any theorem 
% we can replace occurrences of $x$ with $y$ and vice versa.
% Thus, if (hypothetically) $1=2$ and $2=3$, we can replace 
% $2$ with $1$ in the equality $2=3$ to obtain $1=3$.
% Note that $1=3$ is false. 
% But it would indeed be a theorem under the (absurd) hypothesis
% that $1=2$ and $2=3$.
% Thus, we can deduce $2=3 \implies 1=3$ in the context 
% where we have hypothesized $1=2$. 
% And we can deduce $1=2 \implies (2=3 \implies 1=3)$ in the
% general context. This last one is a true theorem,
% valid without additional hypotheses.

Returning to the inference rules of propositional calculus, particularly to
logical negation.
\index{logical negation}%
\mymargin{removal negation}%
For removal, the rule is that of double negation: if $\lnot (\lnot (P))$ is a
theorem, then $P$ is also a theorem. For example, if $\lnot (\lnot (1=1))$ is a
theorem, then $1=1$ is also a theorem.
\mymargin{introduction negation}%
The rule of introduction of logical negation is proof by contradiction.
\index{proof!by contradiction}%
\index{contradiction!proof by}%
If we introduce a formal context in which we assume (for contradiction) that $P$
is a theorem and, in the context, we manage to obtain a contradiction (e.g.,
$\lnot (P)$), then we can assert that (outside the context) $\lnot (P)$ is a
theorem.

For example, if $1=1$ is a theorem, we can prove that $\lnot \lnot (1=1)$ is
also a theorem. Indeed, suppose for contradiction $\lnot \lnot \lnot (1=1)$.
Then, through the elimination of double negation, we can deduce $\lnot (1=1)$.
But this contradicts $\lnot \lnot (1=1)$. Thus, we have proven that $\lnot \lnot
(1=1)$.

The formal rules we have introduced are sufficient to prove the validity of the
truth tables of logical operators.

\begin{exercise}
Try to use the formal inference rules of propositional calculus to prove the
validity of the truth tables.
\end{exercise}

\subsection{predicates, quantifiers}

A \emph{predicate}%
\mymargin{predicate}%
\index{predicate} is a proposition that contains one or more variables, and thus its truth value may depend on the value assigned to the variables. An example of a predicate is $x+2=5$, which is true if $x=3$ and false otherwise.

The variables $x,y,\dots$ are nothing more than symbols of the formal system.
Usually, the last letters of the alphabet are used. Since the letters of the
alphabet are finite in number and since we want to have the possibility of
introducing new variables indefinitely, we decide that a variable can be
composed of a sequence of multiple symbols. For example, one can use the prime:
$x$, $x'$, $x''$ or use a sequence of letters: $\textrm{area}$, $\sin$ to have
an unlimited amount of different names available.

If a predicate depends on one or more variables, these are called \emph{free
variables}%
\mymargin{free variables}%
\index{free variables}. It is possible to \emph{bind} a free variable using a quantifier. The \emph{universal quantifier}%
\mymargin{universal quantifier}%
\index{universal quantifier} denoted by the symbol $\forall$ is used to assert that the predicate is true for every possible value of the quantified variable. For example, the proposition $\forall x\colon x+2=5$ means: ``for every $x$, $x+2=5$ holds'' and is a false proposition. The quantifier $\forall x$ makes the variable $x$ of the predicate $x+2=5$ \emph{mute} in the sense that the truth value of the proposition no longer depends on the value of that variable, which is no longer a free variable but serves only as a placeholder. Indeed, if I change the name of a mute variable, the truth value does not change. For example, writing $\forall x\colon x+1=1+x$ is equivalent to $\forall z\colon z+1=1+z$.

The \emph{existential quantifier}%
\mymargin{existential quantifier}%
\index{existential quantifier} denoted by the symbol $\exists$ is used to assert that the predicate is true for at least one value of the quantified variable. For example, the proposition $\exists x\colon x+2=5$ means: ``there exists at least one $x$ for which $x+2=5$ holds''. What is obtained is a proposition in which the variable $x$ is mute. In this case, the proposition is true because for $x=3$, the predicate is true.

A predicate can depend on multiple variables, and it is therefore possible to
insert multiple quantifiers. In such a case, the order of the quantifiers can be
relevant. Of the following propositions, the first is true, while the second is
false:
\begin{gather*}
\forall x\colon \exists y\colon x+2=y \\
\exists y\colon \forall x\colon x+2=y.
\end{gather*}
Intuitively, the first statement tells us that for every number $x$, there
exists a number $y$ that differs by $2$ from $x$. The second tells us that there
is a number $y$ that differs by $2$ from any number $x$.

\begin{table}
  \begin{center}
    \begin{tabular}{lcl}
    $\lnot \forall x \colon P(x)$ & $\iff$ & $\exists x \colon \lnot P(x)$\\
    $\lnot \exists x \colon P(x)$ & $\iff$ & $\forall x \colon \lnot P(x)$\\
    \end{tabular}
  \end{center}
  \caption{De Morgan's laws for logical quantifiers.}
  \label{tab:proprieta_quantificatori}
\end{table}

It is useful to know how to transform one logical quantifier into another
through negations using De Morgan's laws reported in
table~\ref{tab:proprieta_quantificatori}.

Often, the universal quantifier is implied. For example, if one writes
$x+y=y+x$, it is understood that all free variables are quantified using the
universal quantifier: $\forall x\colon \forall y\colon x+y=y+x$. The punctuation
mark $\colon$ can be omitted, and quantification can be done simultaneously on
multiple variables. We could therefore write: $\forall x\forall y\colon x+y =
y+x$ or even $\forall x,y\colon x+y = y+x$. Sometimes, if it does not create
ambiguity, quantifiers can be written at the end of the sentence
\[
  x+y=y+x \quad \forall x\forall y.
\]

A proposition cannot contain free variables, otherwise its truth value could
depend on those variables and would not be well-defined.
\mymargin{constants}%
\index{constants}%
However, it can be understood that some variables are \emph{constants}, meaning
they represent a specific object. For example, in the proposition $\forall x:
\pi > -x^2$, it is understood that $\pi$ is a constant, i.e., a specific number.
Even the digits $0$, $1$, $2$ are usually always understood as constants.

Along with predicate calculus, the equality symbol \texttt{=} is usually
introduced. How equality is defined depends on how the objects treated by our
formal system are defined, something we will do for set theory in the next
section. However, regardless of how equality is defined, the following axioms
are required to hold.

\begin{axiom}[extensionality]
If $P$ is a predicate in which the variables $x$ and $y$ appear free, and if
$x=y$, then in $P$, I can substitute one or more occurrences of the variable $y$
with the variable $x$ (and vice versa).
\end{axiom}

\begin{axiom}[equality]
\[
  \forall x\colon x=x.  
\]
\end{axiom}

For predicate calculus, we must also specify the inference rules for the
introduction and elimination of quantifiers.

\mymargin{elimination $\forall$}%
\index{quantifier!universal}%
The elimination of the universal quantifier is very simple. If we have a theorem
of the form $\forall x\colon P(x)$ and if a constant $c$ has been introduced in
the demonstrative context, then we can deduce $P(c)$.
\mymargin{introduction $\forall$}%
Conversely, to introduce the universal quantifier, we must open a demonstrative
context in which we assume that $x$ is a constant about which we have no
information (i.e., $x$ is not used as a constant outside the context). If in the
context we manage to prove a proposition $P(x)$, we can then deduce, outside the
context, the theorem $\forall x\colon P(x)$.

In practice, it will be said: given any $x$, we have proven $P(x)$, so we can
assert that $\forall x\colon P(x)$.

\index{quantifier!existential}%
\mymargin{elimination $\exists$}%
Regarding the existential quantifier, the rule of elimination tells us that if
we have a theorem of the form $\exists x\colon P(x)$, we can introduce a
demonstrative context in which a constant $c$ is defined that has the property
$P(c)$. If in that context I manage to obtain a theorem that does not depend on
$c$, then that theorem is valid even outside the context.
\mymargin{introduction $\exists$}%
The rule of introduction tells us instead that if we have a theorem of the form
$P(c)$, with $c$ any constant, we can then deduce $\exists x\colon P(x)$.

As an example, we state and formally prove a simple theorem.

\begin{theorem}[example of formal proof]
  \[ \forall x \exists y\colon x=y \]
\end{theorem}
\begin{proof}
  \begin{align}
    &\forall x\colon x = x & & \text{[axiom]} \label{ff1}\\
    &\{ & &\text{[fix any $x$]}\\
    &\quad x=x  & & \text{[elimination $\forall$ from \eqref{ff1}]}\\ 
    &\quad \exists y\colon x=y & &\text{[introduction $\exists$]}\\
    &\} & &\\
    &\forall x\exists y\colon x=y & &\text{[introduction $\forall$]} 
  \end{align}
\end{proof}